% 對應英文：sections/04_data.tex

\section{面板概覽}

本研究的資料集為一個平衡面板，涵蓋 $T = 324$ 個週度觀測值
（2020-01-05 至 2026-03-15），共 $N = 86$ 種加密貨幣。
原始資料檔保留 53 個特徵欄位；依據最終 v6 實驗設定，
本文在六組資訊集中實際使用其中 42 個模型輸入特徵，
其餘欄位保留於資料檔中供後續擴充，不列入本版實證比較。
橫截面固定為 2026 年 3 月市值排名前 86 的代幣，
尚未上市資產在對應週次的缺失值以 $-99.99$ 標記，
並在損失計算與評估時加以遮罩。

\section{特徵類別}

表~\ref{tab:features} 彙整了本文實際納入模型的 42 個特徵，
可概分為七個群組。
這樣的分組設計反映了本文的核心研究問題：
不同資料型態是否需要不同的模型歸納偏誤，才能轉化為穩健的橫截面排序績效。

\begin{table}[t]
\centering
\caption{本文實際納入模型之特徵類別。$M_g$ 表示第 $g$ 類特徵數。}
\label{tab:features}
\small
\begin{tabular}{llc}
\toprule
\thead{類別} & \thead{代表特徵} & \thead{$M_g$} \\
\midrule
價格與動量         & 1/4/12/26/52 週報酬                         & 5  \\
技術指標           & RSI、布林帶、ATR、OBV、成交量比率          & 6  \\
鏈上指標           & 活躍地址數、交易數、NVT、交易所凈流量、MVRV & 5  \\
總體與市場情緒     & 恐懼貪婪指數、S\&P 500、DXY、VIX、金銀報酬  & 11 \\
ETF 與 Polymarket  & BTC/ETH ETF 流量、Polymarket BTC 機率      & 6  \\
川普社群媒體       & 發文數、大寫比例、關稅分數、加密分數、情緒  & 5  \\
Farside ETF 細粒度 & GBTC、BTC ex-GBTC、ETHE、ETH ex-ETHE       & 4  \\
\midrule
合計               &                                              & 42 \\
\bottomrule
\end{tabular}
\end{table}

價格與技術指標提供最直接的市場行為訊號，
鏈上變數則反映代幣網路的使用與估值狀態；
總體變數捕捉美元流動性與風險偏好的共同環境，
而社群媒體與 ETF 資金流特徵則代表具有事件性與制度性特徵的替代資料。
這種由「慢變基本面」到「快變事件訊號」的結構差異，
也是本文後續比較橫截面模型與時序模型的重要基礎。

\section{資料切割}
\label{sec:data:splits}

我們按時間順序將面板切割為訓練集（70\%，第 0--226 週）、
驗證集（15\%，第 227--274 週）與測試集（15\%，第 275--323 週，共 49 週）。
所有超參數均在驗證集上選擇，測試集僅用於最終一次評估，
以避免在模型設計過程中反覆使用測試期資訊。

\section{前瞻偏差分析}
\label{sec:data:lookahead}

本文特別檢查三種可能的前瞻偏差來源。
第一，所有特徵均以 52 週滾動 Z 分數進行標準化；
在第 $t$ 週，每個特徵只使用區間 $[t-51, t]$ 的均值與標準差，
不會引用未來週次資料。
這種作法與 \citet{gu2020empirical} 的處理精神一致，
即使包含當期值，其在 52 週窗口中的權重也僅約為 $1/52$，
屬金融時間序列研究中的常見作法。

第二，川普社群媒體特徵是以每週結束前可觀察到的貼文內容所構建：
發文數、大寫比例、關稅分數、加密分數與情緒指標，
均僅使用該週內已發布的 Truth Social 與 X 貼文。
第三，ETF 資金流特徵由每日資料聚合為週頻後，再於當週做滾動標準化，
因此亦不包含未來資訊。
綜合而言，本文的主要風險不在於 look-ahead bias，
而在於資料覆蓋率與市場結構變動所造成的估計不穩定。

\section{存活偏差}
\label{sec:data:survival}

本文的 86 種資產係依 2026 年 3 月的市值排名固定選定，
因此不可避免地帶有前瞻性的存活偏差：
若某些 2020 年至 2026 年間曾活躍、但後續衰退或下架的代幣未被納入，
則樣本中的平均報酬與流動性條件可能較實際市場更為樂觀。
這個限制會影響絕對績效水準的解讀，
特別是多空策略在小型、脆弱代幣上的可交易性評估。

然而，存活偏差對本文的跨模型比較不致構成致命問題。
原因在於所有模型皆面對相同的資產宇宙、相同的切割方式與相同的資訊集，
因此比較結果仍可用來判斷不同損失函數與模型架構的相對優劣。
換言之，本文的結論應更被解讀為
「在給定樣本宇宙下，何種方法較適合處理加密貨幣橫截面排序問題」，
而非對整體加密貨幣市場可投資性做無條件外推。
